\chapter{زمان‌بندی پروژه (نمودار گانت)}
\label{app:gantt}

این پیوست زمان‌بندی پروژه را به‌صورت نمای گانت ارائه می‌کند.

%==============================================================================
\section{زمان‌بندی \pnum{12} ماهه}
%==============================================================================

پروژه طبق برنامه‌ی \pnum{12} ماهه‌ی هم‌راستا با پیشنهاد پیش می‌رود. فازهای \pnum{1} تا \pnum{9} مرور ادبیات، طراحی کوهورت، توکن‌سازی، پیاده‌سازی \lr{ETHOS}، آموزش کم‌نمونه و فدرال، مقایسه با خط پایه، آزمایش‌ها و نگارش پایان‌نامه را پوشش می‌دهند.

%==============================================================================
\section{نمای کلی فازها}
%==============================================================================

\begin{table}[h]
\centering
\caption{نمای کلی گانت پروژه}
\begin{tabular}{rrcc}
\toprule
فاز & وظیفه & شروع (هفته) & مدت \\
\midrule
\cnum{1} & مرور ادبیات & \cnum{1} & \cnum{4} \\
\cnum{2} & دسترسی \lr{MIMIC\text{-}IV} و طراحی کوهورت & \cnum{5} & \cnum{2} \\
\cnum{3} & استخراج داده و توکن‌سازی & \cnum{7} & \cnum{4} \\
\cnum{4} & پیاده‌سازی رمزگذار \lr{ETHOS} & \cnum{11} & \cnum{3} \\
\cnum{5} & فرا یادگیری کم‌نمونه & \cnum{14} & \cnum{4} \\
\cnum{6} & آموزش فدرال & \cnum{18} & \cnum{2} \\
\cnum{7} & خط پایه‌ی \lr{Tabular SOTA} & \cnum{12} & \cnum{4} \\
\cnum{8} & آزمایش‌ها و حذف اجزا & \cnum{20} & \cnum{4} \\
\cnum{9} & نگارش و اشکال‌ها & \cnum{24} & \cnum{8} \\
\bottomrule
\end{tabular}
\label{tab:gantt}
\end{table}

%==============================================================================
\section{نقاط عطف کلیدی}
%==============================================================================

\begin{itemize}
\item \textbf{هفته‌ی \pnum{4}}: اتمام مرور ادبیات؛ نهایی‌شدن پرسش‌های پژوهشی.
\item \textbf{هفته‌ی \pnum{7}}: استخراج کوهورت؛ پردازش \pnum{2,000} بیمار اول.
\item \textbf{هفته‌ی \pnum{14}}: عملیاتی‌شدن خط لوله‌ی \lr{ETHOS} و یادگیری کم‌نمونه.
\item \textbf{هفته‌ی \pnum{18}}: اتمام مقایسه‌ی فدرال در برابر متمرکز.
\item \textbf{هفته‌ی \pnum{24}}: اجرای همه‌ی آزمایش‌ها؛ شناسایی بهترین مدل (\lr{RF}).
\item \textbf{هفته‌ی \pnum{32}}: تکمیل پیش‌نویس پایان‌نامه؛ آماده‌سازی دفاع.
\end{itemize}

%==============================================================================
\section{گرافیک نقشه‌ی راه}
%==============================================================================

شکل~\ref{fig:roadmap} نقشه‌ی راه کارهای آینده (\pnum{1} تا \pnum{5} سال) را نشان می‌دهد.

\begin{figure}[htbp]
\centering
\fbox{\parbox{0.8\textwidth}{
\centering
\textbf{نقشه‌ی راه کارهای آینده}\\[0.5em]
فاز \pnum{1} (\pnum{1} سال): کوهورت بزرگ‌تر، اعتبارسنجی بیرونی\\
فاز \pnum{2} (\pnum{2} سال): آزمایش چندسایتی فدرال\\
فاز \pnum{3} (\pnum{3} سال): استقرار آینده‌نگر\\
فاز \pnum{4} (\pnum{4} سال): مسیر تأیید مقرراتی\\
فاز \pnum{5} (\pnum{5} سال): یکپارچه‌سازی پشتیبان تصمیم بالینی
}}
\caption{فازهای کار آینده (\pnum{1} تا \pnum{5} سال)}
\label{fig:roadmap}
\end{figure}

%==============================================================================
\section{مدیریت ریسک}
%==============================================================================

جدول~\ref{tab:risk} ریسک‌های پروژه و راهکارهای کاهش را، مطابق چارچوب مدیریت پروژه‌ی پیشنهاد، خلاصه می‌کند.

\begin{table}[htbp]
\centering
\caption{مدیریت ریسک پروژه}
\footnotesize
\setlength{\tabcolsep}{1.8pt}
\renewcommand{\arraystretch}{1.12}
\begin{tabularx}{\linewidth}{@{}Y C{1.4cm} C{1.4cm} Y@{}}
\toprule
ریسک & احتمال & اثر & کاهش اثر \\
\midrule
تأخیر دسترسی به داده & متوسط & بالا & درخواست زودهنگام اعتبار \lr{MIMIC\text{-}IV}؛ داده‌ی مصنوعی پشتیبان \\
عملکرد ضعیف مدل & پایین & پایین & تکیه بر سقف‌های کلاسیک موجود؛ انباشت کم‌نمونه و خط پایۀ جدولی به‌عنوان مسیر پشتیبان \\
مشکل همگرایی فدرال & پایین & متوسط & تنظیم \lr{FedAvg}؛ بازگشت به متمرکز \\
محدودیت منابع محاسباتی & پایین & پایین & اعتبار ابری؛ کوهورت کوچک‌تر در صورت نیاز \\
لغزش زمان‌بندی & متوسط & متوسط & بافر در گانت؛ اولویت آزمایش‌های هسته \\
\bottomrule
\end{tabularx}
\label{tab:risk}
\end{table}
